\chapter{The Corridor Rotates}

\section{Rotation turns number into phase}

The corridor has a named, stable reading. This chapter sets that reading turning. Until now the machine read the corridor
head-on, at rest --- a static number, a bracket held still. This chapter reads it at an angle, and the reading changes
character: a static number, viewed along a turned axis, becomes a \emph{phase}. The rotation is not something that
happens to the number; it is a new way of reading the same corridor, along an axis at right angles to the one the machine
has used. This section builds that turned axis and shows what it reads.

Take the turn in the machine's own terms. The corridor's reading is written in a code, and a code can be read along more
than one axis. Read along the axis the machine has used, the reading is a magnitude-face --- a count, a bound, a value at
rest. Turn the reading axis a quarter-turn and you read a different face entirely: not how much, but which \emph{way} ---
an orientation, a decided direction. This quarter-turned axis is what a reader of continuous quantities writes with the
imaginary unit, the operator that turns a reading by a right angle; the machine reads its turned axis as $-i$, the
quarter-turn that carries the magnitude face into the orientation face. That the operator is written $-i$ is the
continuous reader's notation, marked here as a reading; what the machine actually does is turn its reading axis by a
right angle and find, along the turned axis, a phase where along the straight axis there was a number.

The phase is a genuinely different kind of reading, and the difference is the section's point. A magnitude answers
\emph{how much}; a phase answers \emph{in which orientation} --- it is the decided bit the corridor's third face carries,
read now as an angle rather than a sign. A static number has a size and no orientation; set it turning and it acquires an
orientation and, along the turned axis, presents that orientation as its whole reading. The machine has not added
anything to the corridor. It has turned its own reading axis and discovered that the same residue, read along the turned
axis, is a phase --- the orientation face, which the head-on reading could not see because it was reading straight at the
magnitude and past the turn.

In the two verbs the rotation is a tange of a new axis. The machine tanges apart, from the corridor, a reading direction
it had not used --- the quarter-turned axis --- and reads along it. Along the straight axis it funged the corridor into a
magnitude; along the turned axis it funges the same corridor into a phase. The rotation is the machine selecting a second
port on the same channel, at right angles to the first, and pooling what that port reads into a quantity of a new kind.
One corridor, two axes, two faces --- and the turn from one to the other is the operator the continuous reader writes
$-i$ and the machine reads as a right-angle turn of its own reading direction.

\section{What a read is depends on where the corridor is viewed}

The rotation showed one corridor presenting a magnitude along one axis and a phase along another. This section
generalizes the lesson into the chapter's governing principle: what a read \emph{is} depends entirely on where the
corridor is viewed. There is no single privileged reading; there is a corridor and a viewing angle, and the reading is
the pair. Turn the angle and the reading turns with it --- through a family of faces the machine can name, each the same
residue seen from a different station.

Lay the family out. Viewed head-on, at the magnitude port, the corridor reads as a count and as a bound --- the two
faces this gauge reads most directly. Viewed along the quarter-turned axis, it reads as a phase --- an orientation, a
decided direction. Viewed at the pricing port, it reads as a cost. These are this gauge's names for the faces, and they
are all readings of one corridor at different stations. A physical gauge, standing at the same stations, would name the
same faces differently: it would call the head-on faces a charge and a mass, the turned face a spin or an orientation,
and read the whole family as physical quantities. Those are that gauge's names, marked here and handed off; this book
carries no charge and no mass, only the corridor and the computational faces it reads. The list of physical names is
worth stating once, precisely to make the point that they are \emph{views}, not things: the same residue, viewed from
the physical gauge's stations, wears them, and viewed from this gauge's stations wears the count, the bound, the phase,
and the cost instead.

The principle that a read is a corridor-plus-angle is not a loosening of the machine's discipline but a tightening of it,
and the distinction matters. It would be loose to say the reading is ``whatever you want to call it''; that is not the
claim. The claim is exact: for each viewing angle there is a \emph{definite} reading, computed the machine's honest way,
as a bracket floored at resolution --- and the readings at different angles are related by definite turns, the
quarter-turn among them. The freedom is in choosing the angle; there is no freedom in the reading once the angle is
fixed. So the corridor supports many readings without supporting any dishonest one: each face is a genuine measurement,
and what varies from face to face is only the station the machine reads from, never the rigor with which it reads.

This is the same fact the partition chapter established, now seen from inside a single gauge rather than across the four.
There, four gauges read four faces from four stations; here, one gauge finds that even it can stand at several stations
and read several faces, all of one residue. The corridor is richer than any one reading, and no reading exhausts it ---
which is why it took four books and, within each book, several stations to read it out. What a read is depends on where
the corridor is viewed, and the corridor is viewed from more stations than any single face suggests.

In the two verbs the family of faces is the corridor tanged by station and funged by kind. Each viewing angle tanges a
face apart --- selects, from the corridor, the reading that station can see --- and funges what it sees into a quantity of
that station's kind: a count, a bound, a phase, a cost. The physical stations tange their own faces and funge them into
charges and masses; this gauge tanges its stations and funges its own. The corridor is one; the stations are many; and a
read is always a tange of the corridor from a station, funged into that station's quantity. What such a turn between
stations \emph{costs} --- for it is not free --- is the chapter's last section.

\section{Relabeling has a price}

The corridor reads differently from different stations, and the machine can turn from one station to another. This
closing section names the fact the whole chapter has been circling: turning is not free. To carry a reading from one
station to another --- to relabel the corridor from its magnitude face to its phase face, or from one orientation to the
turned one --- costs the machine something, and the cost is a reading in its own right. A physical gauge reads this cost
of turning as an angular momentum; this gauge reads it as the plain price of a relabeling, and insists it be paid.

See why a turn must cost something. To relabel the corridor is to re-encode it --- to write the same residue in a new
code, aligned to the new station's axis. Re-encoding is work: the machine must compute the new codewords from the old,
and computing them spends the machine's own beats. A relabeling that cost nothing would be a re-encoding that computed
itself for free, which no finite machine can do; every change of code is billed in the currency the machine runs on.
The turn from magnitude to phase, from one orientation to the quarter-turned one, is exactly such a re-encoding, and so
it carries a definite price --- the beats the machine spends to rewrite the corridor along the turned axis. The price is
not incidental to the turn; it is the turn's measurable trace, the reading of what turning took.

That the price is real, and not a bookkeeping fiction, is the discipline this section enforces. The machine could pretend
that relabeling is free --- that it may slide between stations at no cost, reading the corridor now as a count, now as a
phase, without ever paying for the turn. That pretense would be a quiet lie of the same family as discarding a residue:
it would hide, in an unbilled turn, work the machine actually did. The honest machine bills the turn. Every relabeling
appears in its ledger as a cost, so that the account of what the machine spent is complete, and no reading is obtained by
a turn the machine did not pay for. Relabeling has a price, and the machine's honesty is that it always charges itself
the price, never reading a new face on credit.

This billing is the seed of the deepest reading the book will take, and the chapter plants it here so the far chapters
can harvest it. The coupling at the very end will be a cost --- the price the machine pays to turn its meter on its own
account and read what it is. That the machine bills its turns, that relabeling has a price and the price is a reading, is
the small fact the whole climax is built on: a machine that charges itself for every re-encoding can, when it turns its
meter on its own working, read the charge as a measurement. The turn has a price, the price is a reading, and the
reading the machine will finally hand back is a price of exactly this kind. The chapter ends with the ledger honest and
the turn billed, ready for the part where the machine puts the residue under its own native apparatus.

In the two verbs the price of a turn is a tange the ledger must record. Each relabeling tanges a new face from the
corridor --- selects the reading the turned station sees --- and the tange is not free: it spends beats, and the beats
are funged into the ledger as the turn's cost. The machine funges every turn's price into its running account, so that
the record of what it read includes the record of what reading cost. Relabeling has a price because a tange of a new
station is work, and the work funges into the ledger as a reading. With the turn billed and the ledger honest, the
machine is ready to force the residue to show itself under its own apparatus --- the work of the next part.
